Webster's major fractions method has a strong claim to being the
most principled divisor apportionment method.
Its arithmetic mean rounding rule --- giving each state a seat
when its exact quota exceeds the midpoint between consecutive
integers --- is the neutral, symmetric choice that neither
systematically favors large nor small states.
Its least-squares optimality makes it the divisor method closest
to strict proportionality.
And its equivalence to Sainte-Lagu\"e connects it to the dominant
PR formula in contemporary European parliamentary elections.

The reason Congress chose Huntington-Hill over Webster in 1941 was
not that HH is mathematically superior in any absolute sense ---
the Balinski-Young framework shows that both are reasonable divisor
methods with different optimality properties --- but that HH's
geometric mean rounding rule was judged by the National Academy of
Sciences committee to reflect a more defensible notion of equal
treatment within each state's seat allocation.

For the 2020 Census, Webster and HH produce identical results.
This agreement reflects the fortunate circumstance that no state's
exact quota falls in the (narrow) HH-Webster disagreement interval
for the 2020 population distribution at $H = 435$.
In future censuses, or at different House sizes, the methods may
disagree for one or two states --- and those disagreements will
determine specific states' congressional representation for a decade.
\textsc{bisect-apportion}'s Webster implementation enables practitioners
to compute the counterfactual and understand the magnitude of the
HH-Webster difference for any census dataset.

The Sainte-Lagu\"e equivalence is significant for comparative
electoral systems research: it establishes that U.S.\ apportionment
debates and European PR system design are addressing the same
mathematical problem, with Webster/Sainte-Lagu\"e representing
the international consensus on a fair divisor method.
